\chapter{第二个输出}
% 完整版：chapters/16_chemoutput.tex

这台机器有两个输出。快的是肌肉。慢的是身体的化学：大脑控制心脏、血管和腺体，并通过血液把消息同时发给全身。血液是所有总线中最具广播性的一条：一条消息大约一分钟就能送到每一个细胞，而赋予它含义的，照例是接收者的“锁”。

\section*{油门与刹车}

心脏自己会跳，它有自己的节拍器。大脑只是用两根符号相反的导线来微调节奏。\nt{NORA}是油门踏板：让心跳加快，但加快得慢。\nt{ACET}是刹车：让心跳减慢，而且减得快，因为它不需要一长串化学反应。在这里，\nt{ADRE}类型也终于有了用武之地：释放到血液里的肾上腺素也是油门踏板，只不过是一下子为全身踩下去的。

\pencilfig{16}{\pencilart{16}}{心脏由两根符号相反的导线驱动：油门和刹车。}

\section*{带反馈的级联}

许多激素由级联来控制：大脑命令第一个腺体，第一个命令第二个，第二个分泌激素，而激素再告诉大脑：已经够了。这是一个三级的恒温器。工程师知道这种电路的危险：如果回路调得太“灵敏”，它就会摇摆起来——水平不再保持稳定，而是来回振荡。所以这种调节器的精度是有限的：没法让它既稳定又非常精确。

\pencilfig{127}{%
% three identical first-order stages with proportional feedback: secretion = max(0, K (target - level)).
% K = 3 settles at 3/4 of the target; K = 12 (above the stability limit K = 8) keeps oscillating. x = 0.42 t, y = 0.55 level
\foreach \yo/\t in {1.75/{弱回路：平稳，但达不到目标}, 0/{强回路：更接近目标，但来回摆动}}{
  \draw[pen] (0,\yo) -- (8.5,\yo);
  \draw[pcGray, densely dashed, line width=0.5pt] (0,\yo+0.55) -- (8.5,\yo+0.55);
  \node[lbl, anchor=south west] at (2.0,\yo+0.8) {\t};}
\node[lbl, anchor=west] at (8.55,2.3) {目标}; \node[lbl, anchor=west] at (8.55,0.55) {目标};
\begin{scope}[yshift=1.75cm]
\draw[penlite=pcBlue] plot[smooth] coordinates {(0.00,0.00) (0.17,0.01) (0.34,0.08) (0.50,0.19) (0.67,0.33) (0.84,0.46) (1.01,0.55) (1.18,0.59) (1.34,0.58) (1.51,0.55) (1.68,0.49) (1.85,0.43) (2.02,0.37) (2.18,0.34) (2.35,0.33) (2.52,0.34) (2.69,0.37) (2.86,0.40) (3.02,0.43) (3.19,0.44) (3.36,0.45) (3.53,0.45) (3.70,0.44) (3.86,0.42) (4.03,0.41) (4.20,0.40) (4.37,0.39) (4.54,0.39) (4.70,0.40) (4.87,0.40) (5.04,0.41) (5.38,0.42) (5.88,0.42) (6.38,0.41) (7.06,0.41) (8.40,0.41)};
\end{scope}
\draw[penlite=pcRed] plot[smooth] coordinates {(0.00,0.00) (0.17,0.05) (0.34,0.30) (0.50,0.68) (0.67,1.02) (0.84,1.23) (1.01,1.30) (1.18,1.26) (1.34,1.16) (1.51,1.02) (1.68,0.87) (1.85,0.72) (2.02,0.58) (2.18,0.47) (2.35,0.39) (2.52,0.38) (2.69,0.46) (2.86,0.57) (3.02,0.67) (3.19,0.71) (3.36,0.70) (3.53,0.65) (3.70,0.58) (3.86,0.50) (4.03,0.43) (4.20,0.41) (4.37,0.45) (4.54,0.53) (4.70,0.61) (4.87,0.65) (5.04,0.64) (5.21,0.60) (5.38,0.54) (5.54,0.47) (5.71,0.42) (5.88,0.42) (6.05,0.48) (6.22,0.56) (6.38,0.62) (6.55,0.64) (6.72,0.62) (6.89,0.56) (7.06,0.50) (7.22,0.44) (7.39,0.42) (7.56,0.45) (7.73,0.53) (7.90,0.60) (8.06,0.63) (8.23,0.62) (8.40,0.58)};
\node[lbl, anchor=north] at (4.2,-0.05) {时间};
}{三级激素级联的模型。弱反馈让水平平稳下来，却让它明显低于目标。强反馈让它更接近目标，但水平不再稳住，而是来回摆动。}

\section*{脉冲，而不是水平}

有些激素只能以脉冲方式起作用。如果持续给药，接收的腺体会“疲劳”并停止工作；如果每小时来一次短脉冲，它就正常工作。接收者读的不是水平，而是节律——就像讲摩尔斯电码那一章里只听得见变化的突触。

\section*{昼夜}

最后，身体里还有一个周期约为一天的慢速振荡器。它的自然周期与二十四小时略有出入，每天早晨光线都会校准它，就像收音机调准电台。别把它和计算机的时钟发生器混为一谈：它并不为计算的步骤计时。它切换的是整台机器的工作模式——白天与黑夜，清醒与睡眠。

\fullversion{16}
